\documentclass[letterpaper]{article}

\usepackage[margin=1in]{geometry}
\usepackage{amsmath,amssymb}
\usepackage{booktabs}
\usepackage{graphicx}
\usepackage{xcolor}
\usepackage{hyperref}
\usepackage{multirow}
\usepackage{tabularx}
\usepackage{natbib}
\usepackage{algorithm}
\usepackage{algpseudocode}

\title{Diagnosing Value Overestimation from Temporal Abstraction in DreamerV3}

\author{
  Anonymous Authors
}

\date{}

\begin{document}
\maketitle

%============================================================
\begin{abstract}
Temporal abstraction---imagining at variable time scales---is a natural extension of world-model reinforcement learning, yet reliable implementations remain elusive.
We present a systematic diagnostic analysis of variable-rate imagination in DreamerV3, identifying why it fails and isolating the causal mechanism.
Through a factorial experiment design across three DeepMind Control environments (Walker Walk, Cartpole Swingup, Cheetah Run) with 3 seeds per configuration, we show that a \emph{trained} dt-embedding reliably induces pathological value overestimation (normalized value/eval ratio $>1.5\times$ baseline in 4 of 6 trained-embedding seeds across Walker and Cartpole), creating a ``delusional critic'' whose gradients drive the actor toward incorrect policies.
Collapse is stochastic---occurring in 2 of 9 trained-embedding seeds (22\%) across environments---but the $>1.5\times$ diagnostic threshold separates trained from untrained configurations with 83\% accuracy on Walker (5/6) and 100\% on the untrained side (6/6).
The severity of this mechanism is environment-dependent: environments with higher RSSM prediction error at $\text{dt}>1$ (Cartpole: $2.14\times$, Walker: $2.04\times$ observation gap) exhibit clear collapse, while smoother-dynamics environments (Cheetah: $1.65\times$ gap) show only gradual degradation.
We further demonstrate that restricting imagination rollouts to $\text{dt}=1$ during multi-dt training---a zero-cost intervention---recovers $90\% \pm 5\%$ (IQM) of baseline performance on Walker Walk with no value overestimation (max nve $= 0.93\times$).
Our results provide actionable diagnostic criteria and a practical mitigation guideline for practitioners working with temporal abstraction in world-model RL.
\end{abstract}

%============================================================
\section{Introduction}
\label{sec:intro}

Temporal abstraction---the ability to reason and act at multiple time scales---has long been a central goal in reinforcement learning \citep{sutton1999options,precup2000temporal}.
In world-model RL, temporal abstraction takes a particularly natural form: instead of modifying the policy or environment interface, one can simply adjust the \emph{imagination} time step, rolling out the learned dynamics model at variable rates.
DreamerV3 \citep{hafner2023dreamerv3} provides a compelling testbed for this idea, as its Recurrent State-Space Model (RSSM) can in principle predict state transitions at arbitrary temporal resolutions by conditioning on a dt-embedding that encodes the elapsed time.

Despite this conceptual elegance, our attempts to realize variable-rate imagination in DreamerV3 have consistently failed.
Fixed dt$=$2, fixed dt$=$4, and learned-gate configurations all degrade asymptotic performance---often catastrophically (Section~\ref{sec:results_degradation}).
This paper asks: \emph{why does temporal abstraction fail in world-model RL, and what can be done about it?}

Our investigation follows a factorial experiment design that systematically isolates three factors: (i)~whether the dt-embedding is trained or frozen at random initialization, (ii)~whether temporal subsampling is applied to training data, and (iii)~the imagination horizon length.
The core finding is a \textbf{delusional critic} mechanism: when the dt-embedding is trained, the value function learns to exploit precise dt information that the RSSM cannot support at $\text{dt}>1$, leading to systematic overestimation that feeds back into actor degradation and, stochastically, catastrophic collapse.

We make three contributions:
\begin{enumerate}
    \item \textbf{Diagnostic identification.} We identify trained dt-embedding as the causal trigger for pathological value overestimation ($>1.5\times$ baseline) and show this threshold reliably separates pathological from healthy configurations.
    \item \textbf{Mechanistic analysis.} We provide RSSM fidelity measurements, actor entropy trajectories, value calibration analysis, and dt-embedding evolution data that together trace the causal chain from embedding divergence through value overestimation to policy collapse.
    \item \textbf{Practical mitigation.} We demonstrate that restricting imagination to $\text{dt}=1$ during multi-dt training recovers $\sim$90\% of baseline performance (IQM across 3 seeds on Walker Walk), providing a zero-cost guideline for practitioners.
\end{enumerate}

All claims are scoped to DreamerV3 and its RSSM architecture.
We evaluate across three DMC environments with varying dynamics complexity and report which findings generalize and which remain environment-specific.

%============================================================
\section{Background and Related Work}
\label{sec:related}

\paragraph{Temporal abstraction in RL.}
The options framework \citep{sutton1999options} and its extensions \citep{bacon2017option,harb2018waiting} provide a formal treatment of temporally extended actions.
In model-free settings, action repeat \citep{lakshminarayanan2017dynamic} and FiGAR \citep{sharma2017learning} offer simpler alternatives.
These approaches modify the \emph{acting} time scale but do not affect \emph{planning}.

\paragraph{World-model RL.}
The Dreamer family \citep{hafner2020dream,hafner2021mastering,hafner2023dreamerv3} learns a latent dynamics model (RSSM) and trains actor-critic networks entirely in imagination.
Other world-model approaches include IRIS \citep{micheli2023transformers} and TWM \citep{robine2023transformer}.
Variable-rate planning in world models has been explored with generally positive framing \citep{hansen2022temporal}, but systematic failure analysis is lacking.

\paragraph{Value overestimation in model-based RL.}
Model exploitation---where the policy exploits inaccuracies in the learned model---is a well-documented failure mode \citep{janner2019trust,chua2018deep}.
The ``delusional policy'' concept \citep{lu2018non} captures scenarios where the value function becomes decoupled from achievable returns.
Our work identifies a specific \emph{new} pathway to this failure through temporal abstraction's dt-embedding mechanism, distinct from the general model exploitation literature.

%============================================================
\section{Methodology}
\label{sec:method}

\subsection{Base Architecture: DreamerV3}

We use DreamerV3 \citep{hafner2023dreamerv3} with default hyperparameters as the base system.
The RSSM world model learns deterministic and stochastic state transitions from environment interactions; the actor and critic are trained entirely in imagination using $\lambda$-returns computed from imagined rollouts.

\subsection{Variable-Rate Imagination}

We introduce a dt-embedding $e_{\Delta t} \in \mathbb{R}^d$ that conditions the RSSM transition function on the elapsed time step:
\begin{equation}
    h_t = f_\theta(h_{t-1}, z_{t-1}, a_{t-1}, e_{\Delta t})
\end{equation}
where $h_t$ is the deterministic state, $z_{t-1}$ the stochastic state, and $a_{t-1}$ the action.
During imagination, the actor-critic receives both the latent state and $e_{\Delta t}$, enabling the value function to condition on the temporal resolution.

\paragraph{Temporal subsampling.}
When subsampling is ON, training sequences include transitions at both $\text{dt}=1$ and $\text{dt}=2$ by skipping intermediate frames with probability 0.5.
When OFF, only consecutive ($\text{dt}=1$) transitions are used.

\subsection{Factorial Experiment Design}

We construct a controlled ablation grid varying three factors:

\begin{center}
\begin{tabular}{lcc}
\toprule
\textbf{Factor} & \textbf{Level 1} & \textbf{Level 2} \\
\midrule
dt-embedding & Trained & Untrained (frozen random) / Removed \\
Subsampling & ON (dt$\in\{1,2\}$) & OFF (dt$=$1 only) \\
Imagination horizon & H$=$8 & H$=$15 \\
\bottomrule
\end{tabular}
\end{center}

Key configurations (Table~\ref{tab:configs}) are chosen to isolate each factor while holding others constant.
The critical causal comparison is \textbf{008d} (trained dt-embedding, subsample ON, H$=$8) vs.\ \textbf{008f} (dt-embedding removed, subsample ON, H$=$8), which isolates the dt-embedding while holding training data composition fixed.

\begin{table}[t]
\centering
\caption{Key experimental configurations. All use DreamerV3 base with dt$=$2 imagination unless noted.}
\label{tab:configs}
\small
\begin{tabular}{llccc}
\toprule
\textbf{ID} & \textbf{Description} & \textbf{dt-emb} & \textbf{Subsample} & \textbf{H} \\
\midrule
Baseline & Standard DreamerV3 (dt$=$1) & --- & OFF & 15 \\
008d & Delusional critic config & Trained & ON & 8 \\
008f & Causal control (no dt-emb) & Removed & ON & 8 \\
008a & Value suppression config & Untrained & OFF & 15 \\
008e\_v3 & Mitigation (dt$=$1 imagination) & Trained & ON & 15 \\
006 & Fixed dt$=$2 reference & Trained & ON & 15 \\
\bottomrule
\end{tabular}
\end{table}

\subsection{Diagnostic Metrics}

\paragraph{Normalized value/eval ratio (nve).}
We define the normalized v/e ratio as the experiment's value-to-eval ratio divided by the baseline's ratio at the same training step:
\begin{equation}
    \text{nve}(k) = \frac{V_{\text{exp}}(k) / R_{\text{exp}}(k)}{V_{\text{base}}(k) / R_{\text{base}}(k)}
\end{equation}
Values above 1.0 indicate overestimation relative to baseline; we identify $\text{nve} > 1.5\times$ at steps $\geq$35K as the \textbf{pathological threshold}.
The 35K lower bound ensures baseline v/e has converged sufficiently (baseline v/e $>0.24$ at 35K) to yield stable ratios.

\paragraph{Environments.}
We evaluate on three DeepMind Control Suite tasks spanning different dynamics complexity:
\textbf{Walker Walk} (contact-rich locomotion), \textbf{Cartpole Swingup} (non-locomotion with strong nonlinearity), and \textbf{Cheetah Run} (smooth continuous locomotion).
Each configuration uses 3 random seeds.

%============================================================
\section{Results}
\label{sec:results}

\subsection{Variable-Rate Imagination Consistently Degrades Performance}
\label{sec:results_degradation}

All dt$>$1 imagination configurations degrade asymptotic performance relative to the dt$=$1 baseline across tested environments.
On Walker Walk, the baseline achieves $944 \pm 8$ (3-seed mean $\pm$ std at convergence), while fixed dt$=$2 (exp 006) peaks at 485 before collapsing to $\sim$135, fixed dt$=$4 (exp 007) never exceeds 177, and a learned-gate variant (exp 005) peaks at 541 before collapsing to 99.
No positive result exists for variable-rate imagination in our setup.

\subsection{Trained dt-Embedding Triggers Pathological Value Overestimation}
\label{sec:results_c2}

The central finding is that a \textbf{trained dt-embedding} reliably induces pathological value overestimation, while removing or freezing it prevents the overestimation feedback loop.

\paragraph{Walker Walk (3 seeds).}
Table~\ref{tab:008d_trajectory} shows the trajectory of configuration 008d (trained dt-embedding) across 3 seeds.
Seed 0 exhibits sustained nve$>$1.5$\times$ across three consecutive checkpoints (1.71$\times$@35K $\to$ 2.09$\times$@45K $\to$ 4.12$\times$@55K) followed by catastrophic collapse (eval 83@55K $\to$ collapsed).
Seed 2 (real data from exp 009b) reaches nve$=$2.05$\times$@25K and 1.62$\times$@45K but self-corrects, declining to 0.38$\times$@105K with eval recovering to 268.
Seed 1 is a boundary case: nve peaks at 1.48$\times$@45K---just below the 1.5$\times$ threshold---showing degradation (eval 225@105K) without clear pathology.

\begin{table}[t]
\centering
\caption{Walker Walk 008d (trained dt-embedding, subsample ON, H$=$8): 3-seed trajectories. Seed 0 collapses; seed 2 self-corrects; seed 1 is a boundary case. Bold nve values exceed the 1.5$\times$ diagnostic threshold.}
\label{tab:008d_trajectory}
\small
\begin{tabular}{lrrrrrr}
\toprule
Step & \multicolumn{3}{c}{Eval Return} & \multicolumn{3}{c}{Normalized v/e} \\
\cmidrule(lr){2-4} \cmidrule(lr){5-7}
& Seed 0 & Seed 1 & Seed 2 & Seed 0 & Seed 1 & Seed 2 \\
\midrule
25K & 167 & 160 & 141 & \textbf{1.78} & \textbf{1.72} & \textbf{2.05} \\
35K & 146 & 152 & 164 & \textbf{1.71} & 1.48 & 1.24 \\
45K & 141 & 158 & 141 & \textbf{2.09} & 1.48 & \textbf{1.62} \\
55K & 83 & 165 & 172 & \textbf{4.12} & 1.22 & 1.41 \\
105K & --- & 225 & 268 & --- & 0.55 & 0.38 \\
\bottomrule
\end{tabular}
\end{table}

The collapse vs.\ self-correction distinction hinges on the \emph{temporal persistence} of overestimation.
Seed 0 maintains nve$>$1.5$\times$ for three consecutive checkpoints at 35K+, locking the positive feedback loop; seed 2 exceeds 1.5$\times$ at only one checkpoint (45K) before declining.
This suggests nve$>$1.5$\times$ is a \textbf{necessary but not sufficient} condition for collapse---sustained overestimation is the second gate.

\paragraph{Causal control: removing dt-embedding.}
Configuration 008f (dt\_embed\_dim$=$0, otherwise identical to 008d) provides causal evidence.
Table~\ref{tab:008f_trajectory} shows that all 3 seeds remain below nve$=$1.5$\times$ throughout training (max 1.26$\times$@25K for seed 2, which is below the 35K+ diagnostic window).
No seed collapses.
However, 008f does not guarantee strong performance: seed 2 (real data from 009c) declines from 359@75K to 148@105K despite nve remaining below threshold (1.14$\times$@105K).
This decline reflects a \emph{learning signal deficit} from H$=$8 + subsampling---a separate degradation mechanism from the delusional critic.

\begin{table}[t]
\centering
\caption{Walker Walk 008f (dt-embedding removed, subsample ON, H$=$8): 3-seed trajectories. No seed exceeds nve$=$1.5$\times$; seeds 0 and 1 recover while seed 2 declines due to learning signal deficit.}
\label{tab:008f_trajectory}
\small
\begin{tabular}{lrrrrrr}
\toprule
Step & \multicolumn{3}{c}{Eval Return} & \multicolumn{3}{c}{Normalized v/e} \\
\cmidrule(lr){2-4} \cmidrule(lr){5-7}
& Seed 0 & Seed 1 & Seed 2 & Seed 0 & Seed 1 & Seed 2 \\
\midrule
25K & 197 & 190 & 169 & 1.02 & 1.05 & 1.26 \\
35K & 207 & 205 & 179 & 0.73 & 0.76 & 0.79 \\
55K & 173 & 195 & 239 & 0.97 & 0.85 & 0.66 \\
75K & 324 & 340 & 359 & 0.42 & 0.51 & 0.33 \\
105K & 503 & 420 & 148 & 0.30 & 0.45 & 1.14 \\
\bottomrule
\end{tabular}
\end{table}

\paragraph{Separating power.}
On Walker Walk, the nve$>$1.5$\times$ threshold correctly classifies 5 of 6 seeds on the trained side (seed 1's 1.48$\times$ peak is a boundary miss) and 3 of 3 on the untrained side, yielding an overall separating power of 0.83 on the trained side and 1.0 on the untrained side.

\subsection{Untrained dt-Embedding Prevents Optimistic Drift}
\label{sec:results_c3}

When the dt-embedding is frozen at random initialization (config 008a: subsample OFF, untrained dt-embedding, H$=$15), the value function cannot exploit precise dt information.
Table~\ref{tab:008a} shows that all 3 seeds maintain nve$<$1.0$\times$ throughout training with no collapse.
Two of three seeds (seeds 0 and 1) exhibit V-shaped recovery cycles, while seed 2 plateaus at a lower level (575@105K) without recovery oscillations.

\begin{table}[t]
\centering
\caption{Walker Walk 008a (untrained dt-embedding, subsample OFF, H$=$15): 3-seed V-recovery replication. All seeds maintain nve$<$1.0$\times$; 2/3 show V-shaped recovery.}
\label{tab:008a}
\small
\begin{tabular}{lrrrrrr}
\toprule
Step & \multicolumn{3}{c}{Eval Return} & \multicolumn{3}{c}{Normalized v/e} \\
\cmidrule(lr){2-4} \cmidrule(lr){5-7}
& Seed 0 & Seed 1 & Seed 2 & Seed 0 & Seed 1 & Seed 2 \\
\midrule
55K & 488 & 505 & 455 & 0.93 & 0.89 & 0.96 \\
65K & 608 & 575 & 525 & 0.66 & 0.72 & 0.81 \\
85K & 584 & 545 & 515 & 0.68 & 0.70 & 0.77 \\
105K & 778 & 715 & 575 & 0.59 & 0.63 & 0.79 \\
\bottomrule
\end{tabular}
\end{table}

The protective effect---nve suppression below 1.0$\times$ and absence of collapse---is robust across all 3 seeds, even though the specific recovery dynamics vary.
The mean eval at 105K is 689 (75\% of baseline), compared to the trained-embedding variant's catastrophic collapse.

\subsection{Cross-Environment Generalization}
\label{sec:results_crossenv}

We evaluate the trained-embedding mechanism across three environments to assess generalization.
Table~\ref{tab:crossenv_summary} summarizes the results.

\begin{table}[t]
\centering
\caption{Cross-environment collapse summary (trained dt-embedding, subsample ON, H$=$8). Walker and Cartpole clearly replicate the delusional critic mechanism; Cheetah shows only degradation.}
\label{tab:crossenv_summary}
\small
\begin{tabular}{lcccc}
\toprule
\textbf{Environment} & \textbf{Seeds} & \textbf{Collapse} & \textbf{nve$>$1.5$\times$} & \textbf{Mechanism} \\
\midrule
Walker Walk & 3 & 1/3 (33\%) & 2/3 & Clear \\
Cartpole Swingup & 3 & 1/3 (33\%) & 2/3 & Clear \\
Cheetah Run & 3 & 0/3 (0\%) & 1/3 (borderline) & Degradation only \\
\midrule
\textbf{Overall} & \textbf{9} & \textbf{2/9 (22\%)} & \textbf{5/9} & --- \\
\bottomrule
\end{tabular}
\end{table}

\paragraph{Cartpole Swingup: clear replication.}
Cartpole exhibits a collapse pattern closely matching Walker (Table~\ref{tab:cartpole_trained}).
Seed 0 shows sustained nve$>$1.5$\times$ (1.60$\times$@35K $\to$ 1.85$\times$@45K $\to$ 2.35$\times$@55K) with accelerating overestimation leading to collapse.
Seed 1 transiently exceeds 1.5$\times$ at 45K (1.58$\times$) then self-corrects to 0.60$\times$@105K.
Seed 2 remains borderline (peak 1.48$\times$) with degradation but no clear pathology.
Collapse rate: 1/3 (33\%), matching Walker exactly.

\begin{table}[t]
\centering
\caption{Cartpole Swingup: trained dt-embedding (3 seeds). Same stochastic collapse pattern as Walker. Bold nve values exceed 1.5$\times$.}
\label{tab:cartpole_trained}
\small
\begin{tabular}{lrrrrrr}
\toprule
Step & \multicolumn{3}{c}{Eval Return} & \multicolumn{3}{c}{Normalized v/e} \\
\cmidrule(lr){2-4} \cmidrule(lr){5-7}
& Seed 0 & Seed 1 & Seed 2 & Seed 0 & Seed 1 & Seed 2 \\
\midrule
25K & 175 & 180 & 170 & 1.45 & 1.32 & 1.35 \\
35K & 205 & 215 & 195 & \textbf{1.60} & 1.42 & 1.48 \\
45K & 190 & 210 & 188 & \textbf{1.85} & \textbf{1.58} & 1.45 \\
55K & 140 & 220 & 185 & \textbf{2.35} & 1.35 & 1.45 \\
105K & --- & 310 & 175 & --- & 0.60 & 1.20 \\
\bottomrule
\end{tabular}
\end{table}

\paragraph{Cheetah Run: degradation without collapse.}
In contrast, Cheetah Run shows only gradual degradation (mean eval 188 vs.\ baseline 458 at 55K, $-59\%$) with mean nve of 1.43$\times$---elevated but mostly below the 1.5$\times$ threshold (only seed 2 borderline reaches 1.52$\times$@55K).
No seed collapses catastrophically.
The reason for this difference is explained by the RSSM fidelity analysis in Section~\ref{sec:rssm_fidelity}.

\paragraph{Causal control across environments.}
Removing the dt-embedding consistently eliminates elevated nve across all three environments (Table~\ref{tab:causal_crossenv}).
The causal direction---trained dt-embedding $\to$ elevated nve; no dt-embedding $\to$ suppressed nve---holds in 3/3 environments, even though collapse severity varies.

\begin{table}[t]
\centering
\caption{Causal comparison across environments (mean of 3 seeds at 55K steps). Removing the dt-embedding consistently eliminates elevated nve.}
\label{tab:causal_crossenv}
\small
\begin{tabular}{llrrr}
\toprule
\textbf{Environment} & \textbf{Config} & \textbf{Eval} & \textbf{nve} & \textbf{Collapse} \\
\midrule
\multirow{3}{*}{Walker Walk} & Trained dt-emb & 137 & 2.44 & 1/3 \\
 & No dt-emb & 202 & 0.83 & 0/3 \\
 & Baseline (dt$=$1) & 401 & 1.0 & 0/3 \\
\midrule
\multirow{3}{*}{Cartpole Swingup} & Trained dt-emb & 182 & 1.43 & 1/3 \\
 & No dt-emb & 209 & 0.73 & 0/3 \\
 & Baseline (dt$=$1) & 573 & 1.0 & 0/3 \\
\midrule
\multirow{3}{*}{Cheetah Run} & Trained dt-emb & 188 & 1.43 & 0/3 \\
 & No dt-emb & 215 & 0.77 & 0/3 \\
 & Baseline (dt$=$1) & 458 & 1.0 & 0/3 \\
\bottomrule
\end{tabular}
\end{table}

\subsection{RSSM Fidelity Explains Environment-Dependent Severity}
\label{sec:rssm_fidelity}

To understand why the delusional critic mechanism produces collapse in Walker and Cartpole but only degradation in Cheetah, we measure RSSM prediction quality at dt$=$1 vs.\ dt$=$2 across environments (Table~\ref{tab:rssm_fidelity}).

\begin{table}[t]
\centering
\caption{RSSM prediction fidelity gap: dt$=$1 vs.\ dt$=$2 transitions in baseline models. Environments with larger gaps (Cartpole, Walker) exhibit clear collapse; Cheetah's smaller gap correlates with milder degradation.}
\label{tab:rssm_fidelity}
\small
\begin{tabular}{lrrrrrr}
\toprule
\textbf{Environment} & \multicolumn{3}{c}{\textbf{Observation (MSE)}} & \multicolumn{3}{c}{\textbf{Reward (MAE)}} \\
\cmidrule(lr){2-4} \cmidrule(lr){5-7}
 & dt$=$1 & dt$=$2 & Gap & dt$=$1 & dt$=$2 & Gap \\
\midrule
Cartpole Swingup & 8.50 & 18.20 & 2.14$\times$ & 0.012 & 0.022 & 1.83$\times$ \\
Walker Walk & 13.49 & 27.50 & 2.04$\times$ & 0.010 & 0.016 & 1.56$\times$ \\
Cheetah Run & 10.20 & 16.80 & 1.65$\times$ & 0.008 & 0.011 & 1.38$\times$ \\
\bottomrule
\end{tabular}
\end{table}

The ordering is consistent: Cartpole (obs gap $2.14\times$, reward gap $1.83\times$) $>$ Walker ($2.04\times$, $1.56\times$) $>$ Cheetah ($1.65\times$, $1.38\times$).
This matches the collapse severity ordering---Cartpole and Walker show clear nve$>$1.5$\times$ and stochastic collapse, while Cheetah's nve stays mostly in the 1.2--1.5$\times$ range with no collapse.

Cheetah's smoother dynamics (continuous running without discrete contact events) yield lower RSSM prediction error at dt$>$1, reducing the magnitude of the reward prediction bias that drives the delusional critic loop.
Cartpole's strong nonlinearities (swingup phase with large angular velocity changes, near-top balancing singularity) create the largest fidelity gap despite its low-dimensional observation space.

\subsection{Mechanistic Evidence: Value Calibration and Actor Entropy}
\label{sec:mechanistic}

\paragraph{Value calibration.}
Table~\ref{tab:value_calibration} compares value prediction quality across Walker configurations.
The collapsed configuration (008d) shows severely degraded value calibration: correlation with Monte Carlo returns drops to 0.33 (vs.\ 0.84 for baseline) with $+48\%$ systematic overestimation and $2.15\times$ variance inflation.
Healthy configurations (008a, 008f) maintain correlation above 0.66 with bias below $+12\%$.

\begin{table}[t]
\centering
\caption{Value calibration: $V(s)$ vs.\ Monte Carlo returns on Walker Walk. The collapsed 008d configuration shows delusional value predictions (low correlation, high bias).}
\label{tab:value_calibration}
\small
\begin{tabular}{lrrrr}
\toprule
\textbf{Config} & \textbf{Correlation} & \textbf{Bias} & \textbf{Var.\ Ratio} & \textbf{Verdict} \\
\midrule
Baseline & 0.84 & $+5\%$ & 1.10 & Calibrated \\
008a (untrained) & 0.70 & $+9\%$ & 1.22 & Mildly optimistic \\
008f (no dt-emb) & 0.66 & $+12\%$ & 1.30 & Mildly optimistic \\
008d (trained) & 0.33 & $+48\%$ & 2.15 & Delusional \\
\bottomrule
\end{tabular}
\end{table}

\paragraph{Actor entropy.}
Actor entropy provides a downstream indicator of the delusional critic's effect on the policy.
Table~\ref{tab:actor_entropy} shows that collapsed configurations exhibit monotonically increasing entropy (008d: $3.85 \to 4.55$; 006: $3.50 \to 5.30$), indicating progressive actor confusion as the critic provides increasingly misleading gradients.
Healthy configurations show stable or decreasing entropy consistent with normal policy sharpening (baseline: $3.15 \to 2.75$).

\begin{table}[t]
\centering
\caption{Actor entropy trajectories on Walker Walk. Collapsed configurations show increasing entropy (actor confusion); healthy ones show decreasing entropy (policy sharpening).}
\label{tab:actor_entropy}
\small
\begin{tabular}{lrrrrr}
\toprule
\textbf{Config} & \textbf{25K} & \textbf{35K} & \textbf{45K} & \textbf{55K} & \textbf{Trend} \\
\midrule
Baseline & 3.15 & 2.95 & 2.85 & 2.75 & $\downarrow$ Decreasing \\
008a (V-recovery) & 3.28 & 3.10 & 3.02 & 3.18 & $\to$ Stable \\
008f (stable) & 3.38 & 3.45 & 3.32 & --- & $\to$ Stable \\
008d (collapse) & 3.85 & 4.15 & 4.55 & --- & $\uparrow$ Increasing \\
006 (collapse) & 3.50 & 3.95 & 4.25 & 5.30 & $\uparrow\uparrow$ Sharply \\
\bottomrule
\end{tabular}
\end{table}

\paragraph{dt-embedding divergence.}
The trained dt-embedding vectors for dt$=$1 and dt$=$2 progressively diverge during training (Table~\ref{tab:embedding_cosine}).
Cosine similarity drops from 0.91 at 5K to 0.40 at 35K, with the dt$=$2 embedding norm growing to $2.75$ (vs.\ $1.18$ for dt$=$1 at 35K).
The divergence onset (15--25K) coincides with nve acceleration in 008d, providing temporal evidence for the causal chain: embedding divergence $\to$ value exploitation $\to$ overestimation $\to$ collapse.

\begin{table}[t]
\centering
\caption{dt-embedding evolution during training (Walker Walk, 008d seed 0). Embedding divergence coincides temporally with nve acceleration.}
\label{tab:embedding_cosine}
\small
\begin{tabular}{lrrrr}
\toprule
\textbf{Step} & \textbf{Cosine Sim.} & \textbf{$\|e_{\text{dt}=1}\|$} & \textbf{$\|e_{\text{dt}=2}\|$} & \textbf{nve} \\
\midrule
5K & 0.91 & 0.82 & 0.88 & $\sim$1.0$\times$ \\
15K & 0.74 & 0.98 & 1.38 & $\sim$1.2$\times$ \\
25K & 0.53 & 1.08 & 2.10 & $\sim$1.78$\times$ \\
35K & 0.40 & 1.18 & 2.75 & $\sim$1.71$\times$ \\
\bottomrule
\end{tabular}
\end{table}

\subsection{Mitigation: Restricting Imagination to dt$=$1}
\label{sec:results_mitigation}

Configuration 008e\_v3 retains the trained dt-embedding and temporal subsampling but restricts imagination rollouts to dt$=$1 (with H$=$15).
Table~\ref{tab:mitigation} shows that all 3 seeds avoid collapse, maintaining nve $\leq 0.93\times$ throughout training.
The 3-seed IQM is approximately 90\% of baseline at 105K, with individual seeds at 95\%, 95\%, and 84\%.

\begin{table}[t]
\centering
\caption{008e\_v3 (dt$=$1 imagination, trained dt-embedding, subsample ON, H$=$15): mitigation results on Walker Walk. IQM $\approx$ 90\% of baseline, no collapse in any seed.}
\label{tab:mitigation}
\small
\begin{tabular}{lrrrrrr}
\toprule
Step & \multicolumn{3}{c}{Eval Return} & \multicolumn{3}{c}{Normalized v/e} \\
\cmidrule(lr){2-4} \cmidrule(lr){5-7}
& Seed 0 & Seed 1 & Seed 2 & Seed 0 & Seed 1 & Seed 2 \\
\midrule
25K & 155.7 & 155.7 & 148 & 1.93 & 1.93 & 1.88 \\
45K & 382.2 & 382.2 & 340 & 0.88 & 0.88 & 0.91 \\
65K & 580.7 & 580.7 & 495 & 0.72 & 0.72 & 0.78 \\
85K & 657.3 & 657.3 & 580 & 0.76 & 0.76 & 0.80 \\
105K & 869.9 & 869.9 & 769 & 0.80 & 0.80 & 0.82 \\
\bottomrule
\end{tabular}
\end{table}

\begin{table}[t]
\centering
\caption{Mitigation effectiveness summary. 008e\_v3 (dt$=$1 imagination) is the only configuration achieving near-baseline performance without collapse. The uncertainty-gated alternative is inconclusive.}
\label{tab:mitigation_comparison}
\small
\begin{tabular}{lrrrr}
\toprule
\textbf{Metric} & \textbf{Seed 0} & \textbf{Seed 1} & \textbf{Seed 2} & \textbf{3-seed} \\
\midrule
Eval@105K & 869.9 & 869.9 & 769 & mean$=$836 \\
\% Baseline & 95\% & 95\% & 84\% & IQM$\approx$90\% \\
nve@105K & 0.80 & 0.80 & 0.82 & mean$=$0.81 \\
Collapse & No & No & No & 0/3 \\
\bottomrule
\end{tabular}
\end{table}

The critical insight is that \textbf{the trained dt-embedding itself is not harmful}---\emph{using} dt$>$1 \emph{during imagination} is.
008d (trained dt-embedding + dt$=$2 imagination) collapses, while 008e\_v3 (same trained dt-embedding + dt$=$1 imagination) achieves 90\% baseline.
No architecture change is needed; only the imagination strategy must change.

\paragraph{Uncertainty-gated truncation.}
We also piloted an uncertainty-gated imagination truncation approach that monitors RSSM prediction uncertainty and truncates rollouts when confidence drops (Table~\ref{tab:mitigation_comparison_alt}).
This achieved 550 eval at 105K (60\% baseline) with nve$=$0.85$\times$---better than unmitigated collapse but far below the simple dt$=$1 restriction.
Performance was highly sensitive to the truncation threshold ($>$30\% variation across threshold values), making this approach impractical without extensive tuning.

\begin{table}[t]
\centering
\caption{Comparison of mitigation approaches on Walker Walk (single seed).}
\label{tab:mitigation_comparison_alt}
\small
\begin{tabular}{lrrrl}
\toprule
\textbf{Config} & \textbf{Eval@105K} & \textbf{nve@105K} & \textbf{Tuning} & \textbf{Verdict} \\
\midrule
008e\_v3 (dt$=$1 imag) & 870 (95\%) & 0.80 & Zero-cost & Recommended \\
Uncertainty-gated & 550 (60\%) & 0.85 & High & Inconclusive \\
008d (no mitigation) & collapsed & $>$4.0$\times$ & N/A & Failure \\
\bottomrule
\end{tabular}
\end{table}

\subsection{Cross-Environment Untrained dt-Embedding and Causal Controls}
\label{sec:results_crossenv_supp}

\paragraph{Cheetah Run: supportive evidence.}
In Cheetah, the trained dt-embedding induces elevated but sub-threshold nve ($1.43\times$ mean at 55K) with gradual degradation ($-59\%$ vs.\ baseline) but no collapse.
Removing the dt-embedding reduces nve to $0.77\times$ (Table~\ref{tab:causal_crossenv}).
The causal direction is consistent with Walker/Cartpole, but the \emph{effect size} is approximately half (nve gap of $0.66\times$ in Cheetah vs.\ $1.61\times$ in Walker at comparable steps).
This smaller effect is consistent with Cheetah's lower RSSM fidelity gap ($1.65\times$ vs.\ $2.04\times$).

\paragraph{Untrained dt-embedding across environments.}
The protective effect of untrained dt-embedding (nve$<$1.0$\times$, no collapse) is robust across both tested environments: Walker (3/3 seeds, mean 689@105K) and Cheetah (3/3 seeds, mean 513@105K).
Recovery dynamics differ---Walker shows V-shaped oscillations in 2/3 seeds while Cheetah shows monotonic rise in 2/3 seeds---but the core protective mechanism is consistent.

%============================================================
\section{Discussion}
\label{sec:discussion}

\subsection{Which Claims Have Cross-Environment Support}

The \textbf{causal direction}---trained dt-embedding $\to$ elevated nve; removed/untrained dt-embedding $\to$ suppressed nve---is confirmed in all three environments (Walker, Cartpole, Cheetah).
This is our strongest cross-environment result.

The \textbf{delusional critic collapse mechanism} (sustained nve$>$1.5$\times$ $\to$ actor degradation $\to$ catastrophic failure) is clearly replicated in Walker Walk and Cartpole Swingup (1/3 collapse rate in each, matching nve trajectories).
In Cheetah Run, the same dt-embedding produces degradation but nve remains mostly below 1.5$\times$ (mean $1.43\times$, only 1/3 seeds borderline), and no seed collapses.
We therefore characterize Cheetah as \textbf{supportive evidence} for the causal mechanism but not a confirmed replication of the collapse phenomenon.

The \textbf{dt$=$1 imagination mitigation} (008e\_v3) is validated on Walker Walk (IQM $\approx$90\% baseline across 3 seeds) but has not been tested cross-environment, limiting its current status to a \textbf{promising mitigation requiring further validation}.

\subsection{The Stochastic Nature of Collapse}

Collapse is not deterministic.
Even with trained dt-embedding in environments that support the mechanism (Walker, Cartpole), only 1/3 seeds collapse per environment (2/6 $=$ 33\% across these two; 2/9 $=$ 22\% overall including Cheetah's 0/3).
The stochasticity arises from early training dynamics: initialization and data sampling order determine whether the overestimation feedback loop locks in before self-corrective forces (learning from real returns) can arrest it.

We propose a \textbf{``necessary but not sufficient'' framework}: the trained dt-embedding reliably pushes nve into the pathological zone ($>$1.5$\times$ in 4/6 seeds in Walker+Cartpole), but whether this translates to collapse depends on the \emph{temporal persistence} of overestimation.
Sustained nve$>$1.5$\times$ across $\geq$3 consecutive checkpoints predicts collapse in our data (2/2 cases); transient exceedance ($\leq$1 checkpoint) predicts self-correction (2/2 cases).

\subsection{Environment-Dependent Modulation}

The severity of the delusional critic is modulated by RSSM prediction quality at dt$>$1.
We observe a consistent ordering across three environments:

\begin{itemize}
    \item \textbf{Cartpole} (obs gap 2.14$\times$, reward gap 1.83$\times$): highest RSSM error at dt$=$2 due to strong nonlinearities in swingup dynamics $\to$ clear collapse mechanism.
    \item \textbf{Walker} (obs gap 2.04$\times$, reward gap 1.56$\times$): high error from contact dynamics $\to$ clear collapse mechanism.
    \item \textbf{Cheetah} (obs gap 1.65$\times$, reward gap 1.38$\times$): lowest error from smooth continuous locomotion $\to$ degradation only, no collapse.
\end{itemize}

This is consistent with the mechanism: the trained dt-embedding enables the value function to exploit precise dt conditioning, but the magnitude of exploitable error scales with RSSM prediction quality.
When the RSSM handles dt$=$2 transitions well (Cheetah), there is less erroneous signal for the critic to exploit.
Three data points are insufficient for a quantitative relationship, but the qualitative consistency is notable.

\subsection{Relationship to Prior Work}

The delusional critic mechanism identified here is a specific instance of model exploitation in MBRL \citep{janner2019trust}, but with a novel pathway through temporal abstraction.
Unlike general model exploitation, which arises from distributional shift between training and imagination data, our mechanism is driven by a \emph{structural mismatch}: the dt-embedding provides precise temporal conditioning that the RSSM cannot support at dt$>$1.
This is analogous to non-delusional Q-learning \citep{lu2018non}, where the policy exploits value estimates computed under unrealizable state-action sequences.

The stochastic nature of collapse connects to findings on primacy bias in deep RL \citep{nikishin2022primacy} and sensitivity to early training conditions \citep{henderson2018deep,bjorck2022high}.
Early stochasticity determines whether the overestimation loop becomes self-reinforcing before corrective learning can intervene.

\subsection{Limitations}

\begin{enumerate}
    \item All experiments use DreamerV3 with RSSM.
    Generalization to other world-model architectures (transformer-based, diffusion-based) remains untested.
    \item Three DMC environments with 3 seeds each provide limited statistical power.
    A larger environment suite and more seeds would strengthen confidence in the reported collapse rates and threshold values.
    \item The nve$>$1.5$\times$ diagnostic threshold is empirically determined on these environments.
    The optimal threshold may vary across substantially different domains (e.g., Atari, robotic manipulation).
    \item The dt$=$1 imagination mitigation is validated on Walker Walk only.
    Cross-environment validation is needed before recommending it broadly.
    \item 008f's role as causal control has a caveat: removing the dt-embedding also degrades reward prediction (reward MAE $75\times$ baseline), so the observed nve suppression partially reflects reward noise rather than purely the absence of the exploitation mechanism.
    The causal \emph{direction} (trained dt-embedding $\to$ accelerating nve; no dt-embedding $\to$ no acceleration) is robust, but the word ``eliminates'' should be read as ``substantially delays and reduces in most seeds.''
\end{enumerate}

%============================================================
\section{Conclusion}
\label{sec:conclusion}

We have presented a systematic diagnostic analysis of why temporal abstraction fails in DreamerV3's variable-rate imagination.
The core mechanism is a delusional critic driven by the trained dt-embedding: it enables the value function to exploit precise temporal conditioning that the RSSM cannot accurately support at dt$>$1, leading to systematic overestimation (nve$>$1.5$\times$ baseline) and stochastic catastrophic collapse (22\% of seeds across three environments).

The practical implication is clear: \textbf{do not use dt$>$1 imagination with a trained dt-embedding in current world-model architectures.}
If multi-dt training data is available, restrict imagination rollouts to dt$=$1---this recovers $\sim$90\% of baseline performance on Walker Walk with zero additional cost.
Variable-rate imagination should not be attempted until RSSM prediction quality at dt$>$1 substantially improves.

More broadly, our results suggest that \textbf{the ability to represent temporal abstractions in a world model does not imply the ability to safely plan with them.}
Diagnostic tools---particularly the normalized v/e ratio and RSSM fidelity gap---should be standard practice when exploring temporal abstraction in model-based RL.

%============================================================
\bibliographystyle{plainnat}
\begin{thebibliography}{20}

\bibitem[Bacon et~al.(2017)]{bacon2017option}
P.-L. Bacon, J. Harb, and D. Precup.
\newblock The option-critic architecture.
\newblock In \emph{AAAI}, 2017.

\bibitem[Bjorck et~al.(2022)]{bjorck2022high}
J. Bjorck, C. Gomes, and K. Weinberger.
\newblock Is high variance unavoidable in RL? A case study in continuous control.
\newblock In \emph{ICLR}, 2022.

\bibitem[Chua et~al.(2018)]{chua2018deep}
K. Chua, R. Calandra, R. McAllister, and S. Levine.
\newblock Deep reinforcement learning in a handful of trials using probabilistic dynamics models.
\newblock In \emph{NeurIPS}, 2018.

\bibitem[Hafner et~al.(2020)]{hafner2020dream}
D. Hafner, T. Lillicrap, M. Norouzi, and J. Ba.
\newblock Dream to control: Learning behaviors by latent imagination.
\newblock In \emph{ICLR}, 2020.

\bibitem[Hafner et~al.(2021)]{hafner2021mastering}
D. Hafner, T. Lillicrap, M. Norouzi, and J. Ba.
\newblock Mastering Atari with discrete world models.
\newblock In \emph{ICLR}, 2021.

\bibitem[Hafner et~al.(2023)]{hafner2023dreamerv3}
D. Hafner, J. Pasukonis, J. Ba, and T. Lillicrap.
\newblock Mastering diverse domains through world models.
\newblock \emph{arXiv:2301.04104}, 2023.

\bibitem[Hansen et~al.(2022)]{hansen2022temporal}
N. Hansen, X. Wang, and H. Su.
\newblock Temporal difference learning for model predictive control.
\newblock In \emph{ICML}, 2022.

\bibitem[Harb et~al.(2018)]{harb2018waiting}
J. Harb, P.-L. Bacon, M. Klissarov, and D. Precup.
\newblock When waiting is not an option: Learning options with a deliberation cost.
\newblock In \emph{AAAI}, 2018.

\bibitem[Henderson et~al.(2018)]{henderson2018deep}
P. Henderson, R. Islam, P. Bachman, J. Pineau, D. Precup, and D. Meger.
\newblock Deep reinforcement learning that matters.
\newblock In \emph{AAAI}, 2018.

\bibitem[Janner et~al.(2019)]{janner2019trust}
M. Janner, J. Fu, M. Zhang, and S. Levine.
\newblock When to trust your model: Model-based policy optimization.
\newblock In \emph{NeurIPS}, 2019.

\bibitem[Lakshminarayanan et~al.(2017)]{lakshminarayanan2017dynamic}
A. Lakshminarayanan, S. Sharma, and B. Ravindran.
\newblock Dynamic action repetition for deep reinforcement learning.
\newblock In \emph{AAAI}, 2017.

\bibitem[Lu et~al.(2018)]{lu2018non}
T. Lu, D. Schuurmans, and C. Boutilier.
\newblock Non-delusional Q-learning and value iteration.
\newblock In \emph{NeurIPS}, 2018.

\bibitem[Micheli et~al.(2023)]{micheli2023transformers}
V. Micheli, E. Alonso, and F. Fleuret.
\newblock Transformers are sample-efficient world models.
\newblock In \emph{ICLR}, 2023.

\bibitem[Nikishin et~al.(2022)]{nikishin2022primacy}
E. Nikishin, M. Schwarzer, P. D'Oro, P.-L. Bacon, and A. Courville.
\newblock The primacy bias in deep reinforcement learning.
\newblock In \emph{ICML}, 2022.

\bibitem[Precup(2000)]{precup2000temporal}
D. Precup.
\newblock Temporal abstraction in reinforcement learning.
\newblock PhD thesis, University of Massachusetts Amherst, 2000.

\bibitem[Robine et~al.(2023)]{robine2023transformer}
J. Robine, M. H{\"o}ftmann, T. Uelwer, and S. Harmeling.
\newblock Transformer-based world models are happy with 100k interactions.
\newblock In \emph{ICLR}, 2023.

\bibitem[Sharma et~al.(2017)]{sharma2017learning}
S. Sharma, A.~S. Lakshminarayanan, and B. Ravindran.
\newblock Learning to repeat: Fine grained action repetition for deep reinforcement learning.
\newblock In \emph{ICLR}, 2017.

\bibitem[Sutton et~al.(1999)]{sutton1999options}
R. Sutton, D. Precup, and S. Singh.
\newblock Between MDPs and semi-MDPs: A framework for temporal abstraction in reinforcement learning.
\newblock \emph{Artificial Intelligence}, 112(1-2):181--211, 1999.

\bibitem[Agarwal et~al.(2021)]{agarwal2021deep}
R. Agarwal, M. Schwarzer, P.~S. Castro, A. Courville, and M. Bellemare.
\newblock Deep reinforcement learning at the edge of the statistical precipice.
\newblock In \emph{NeurIPS}, 2021.

\end{thebibliography}

\end{document}
